% =====================================================================
%  Module 5 : Projet final
% =====================================================================
\section{Projet final}

\begin{frame}{Objectif}
  \begin{itemize}
    \setlength{\itemsep}{.5em}
    \item Déployer une application web publique sur une instance EC2
    \item Fonctions : déposer un fichier (10 Mo au plus), lister les fichiers, les télécharger
    \item Fichiers stockés dans un bucket S3 privé
    \item Application fournie : Flask + boto3 ; travail demandé : le déploiement
  \end{itemize}
  \bigskip
  \textbf{Exigences}
  \begin{itemize}
    \item HTTP (80) et HTTPS (443) ouverts à tous, SSH (22) limité à votre adresse
    \item Aucune clé d'accès sur l'instance : accès à S3 par un rôle IAM
    \item Rôle limité au bucket : lister, lire, écrire, pas de suppression
    \item Redémarrage automatique de l'application (panne, redémarrage de l'instance)
  \end{itemize}
\end{frame}

\begin{frame}{Architecture}
  \centering
  \includegraphics[width=\textwidth]{projet-architecture}
\end{frame}

\begin{frame}{Composants sur l'instance}
  \renewcommand{\arraystretch}{1.4}
  \begin{tabularx}{\textwidth}{@{}>{\raggedright\arraybackslash}p{3.6cm}>{\raggedright\arraybackslash}X@{}}
    \toprule
    \textbf{Nginx} & serveur frontal sur les ports 80 et 443 : TLS, taille maximale des requêtes, relais vers l'application \\
    \textbf{Gunicorn} & serveur d'application Python, 2 processus, écoute sur \texttt{127.0.0.1:8000} uniquement \\
    \textbf{Flask + boto3} & application ; boto3 obtient les identifiants du rôle via IMDSv2 \\
    \textbf{systemd} & service \texttt{galerie} : démarrage au boot, relance en cas d'arrêt \\
    \textbf{Utilisateur \texttt{galerie}} & compte système sans droits d'administration \\
    \textbf{\texttt{/etc/galerie.env}} & configuration : nom du bucket, région ; aucun secret \\ \bottomrule
  \end{tabularx}
\end{frame}

\begin{frame}{Traitement d'une requête}
  \centering
  \includegraphics[height=.78\textheight]{projet-requete}
\end{frame}

\begin{frame}{URL présignées}
  \begin{itemize}
    \setlength{\itemsep}{.55em}
    \item URL signée avec les identifiants de l'application, pour un objet et une durée donnés
    \item Générée localement, sans appel réseau
    \item Le navigateur télécharge directement depuis S3 : pas de charge sur l'instance
    \item Durée de validité limitée par celle des identifiants : quelques heures avec un rôle, 7 jours au maximum avec une clé permanente
  \end{itemize}
  \lien{https://docs.aws.amazon.com/AmazonS3/latest/userguide/using-presigned-url.html}
\end{frame}

\begin{frame}{HTTPS}
  \begin{itemize}
    \setlength{\itemsep}{.55em}
    \item Certificat reconnu par les navigateurs : délivré pour un nom de domaine
    \item Projet : certificat autosigné, connexion chiffrée mais avertissement du navigateur
    \item En production :
    \begin{itemize}
      \item nom de domaine (Route 53 ou autre registraire)
      \item certificat gratuit AWS Certificate Manager, sur un répartiteur de charge ou CloudFront
      \item ou certificat Let's Encrypt directement sur Nginx
    \end{itemize}
  \end{itemize}
\end{frame}

\begin{frame}{Vérifications attendues}
  \renewcommand{\arraystretch}{1.15}
  \begin{tabularx}{\textwidth}{@{}>{\raggedright\arraybackslash}p{4.6cm}>{\raggedright\arraybackslash}X@{}}
    \toprule
    \textbf{Exigence} & \textbf{Vérification} \\ \midrule
    Application publique & \texttt{http://IP} et \texttt{https://IP} répondent \\
    SSH restreint & port 22 fermé depuis CloudShell \\
    Application non exposée & port 8000 fermé ; \texttt{ss -ltn} montre \texttt{127.0.0.1:8000} \\
    Bucket privé & URL directe d'un objet : \texttt{AccessDenied} ; URL présignée : accès puis expiration \\
    Aucune clé & pas de \texttt{\textasciitilde/.aws}, aucune clé dans l'application \\
    Moindre privilège & \texttt{aws s3 rm} refusé \\
    Résilience & relance après \texttt{kill -9} et après redémarrage \\ \bottomrule
  \end{tabularx}
\end{frame}

\begin{frame}{Limites et évolutions}
  \begin{columns}[c]
    \begin{column}{.57\textwidth}
      \includegraphics[width=\textwidth]{architecture-cible}
    \end{column}
    \begin{column}{.41\textwidth}
      \begin{itemize}
        \setlength{\itemsep}{.4em}
        \item Une instance, une AZ : pas de haute disponibilité
        \item Adresse IP changeante : nom de domaine, Elastic IP ou répartiteur
        \item Évolutions : ALB + groupe Auto Scaling multi-AZ, RDS Multi-AZ, alarmes CloudWatch
        \item Infrastructure décrite en code : CloudFormation, CDK, Terraform
      \end{itemize}
    \end{column}
  \end{columns}
  \lien{https://docs.aws.amazon.com/wellarchitected/latest/framework/welcome.html}
\end{frame}

\diapotp{5}{Déployer la galerie}{%
  \item Archive de l'application dans le bucket
  \item Instance avec rôle, installation, service systemd
  \item Nginx en HTTP puis en HTTPS
  \item Vérifications, puis suppression de toutes les ressources}


\begin{frame}{Pour aller plus loin}
  \begin{itemize}
    \setlength{\itemsep}{.5em}
    \item Documentation AWS : \url{https://docs.aws.amazon.com}
    \item AWS Well-Architected Framework : six piliers (excellence opérationnelle, sécurité, fiabilité, performance, coûts, durabilité)
    \item AWS Skill Builder : formations et laboratoires en ligne
    \item Certification d'entrée : AWS Certified Cloud Practitioner
    \item Architectures de référence : AWS Solutions Library, \url{https://aws.amazon.com/solutions/}
  \end{itemize}
\end{frame}
